\section{Background }
\label{sec:pre}
% \subsection{Preliminary}
% \para{Quantum bits, gates and circuits}
\begin{figure}[t]
  \centering
  % \includegraphics[width=\linewidth]{images/quantum_gates.drawio.pdf}
  \includegraphics[width=0.65\linewidth]{images/quantum_gates.drawio_cropped.pdf}
     % \vspace{-0.4cm}
  \caption{\revisionA{Quantum gates represented as matrices.  } }
  \label{fig:gates}
   % \vspace{-0.4cm}
\end{figure}

\subsection{Quantum Bits, Gates, and Circuits}

A classical bit stores either 0 or 1. In contrast, a \emph{quantum bit} (qubit) can be represented as a unit vector in a two-dimensional complex vector space, typically described using the computational basis of $\ket{0}$ and $\ket{1}$:
$$
\ket{0}=\begin{bmatrix}1\\0\end{bmatrix},
\qquad
\ket{1}=\begin{bmatrix}0\\1\end{bmatrix}.
$$
An arbitrary single-qubit state can be written as a superposition of these basis states:
$$
\ket{\psi}=\alpha\ket{0}+\beta\ket{1},
\qquad
|\alpha|^2+|\beta|^2=1,
$$
where $\alpha$ and $\beta$ are complex amplitudes. 
If we perform a \emph{measurement} on the quantum state $\ket{\psi}$, it collapses to a classical bit of 0 or 1. That is, a measurement of $\ket{\psi}$ against the computational basis yields outcome $0$ with probability $|\alpha|^2$ and outcome $1$ with probability $|\beta|^2$. 
% Measurements in other bases can be realized by first applying a suitable unitary transformation followed by a computational-basis measurement. Since measurement reveals only a single sampled outcome, quantum computation does not provide direct access to exponentially many parallel results.
An $N$-qubit state can be expressed as a superposition over computational basis states indexed by bitstrings:
$$
\ket{\Psi}=\sum_{x\in\{0,1\}^N} c_x \ket{x},
\qquad
\sum_{x\in\{0,1\}^N} |c_x|^2=1.
$$
Fix an ordering of the computational basis states (e.g., lexicographic order). The \emph{state vector}
representation of $\ket{\Psi}$ is the column vector of its amplitudes in that order:
$
\left[\begin{smallmatrix}
c_{0\cdots 0}\\
c_{0\cdots 1}\\
\vdots\\
c_{1\cdots 1}
\end{smallmatrix}\right].
$
For example, when $N=2$, any two-qubit state $\ket{\Psi}$ can be represented as a linear combination of
$\ket{00}$, $\ket{01}$, $\ket{10}$, and $\ket{11}$, and
$|c_{00}|^2+|c_{01}|^2+|c_{10}|^2+|c_{11}|^2=1$.
The state vector representation of $\ket{\Psi}$ is
$
\left[\begin{smallmatrix}
c_{00}\\
c_{01}\\
c_{10}\\
c_{11}
\end{smallmatrix}\right].
$


 

\begin{figure}[t]
  \centering
  % \includegraphics[width=\linewidth]{images/Quantum circuit background.png}
  \begin{quantikz}
        \lstick{$\ket{0}$} & \gate{H} & \ctrl{1} &  &\\
        \lstick{$\ket{0}$} &  &\targ{} & \ctrl{1} & \\
        \lstick{$\ket{0}$} &  & &\targ{} & \\
    \end{quantikz}
     \vspace{-0.4cm}
  \caption{GHZ state preparation circuit: it takes the all-zero state  $\ket{000}$ as input and prepares a maximally entangled state.}
 \vspace{-0.4cm}
  \label{fig:subcircuits}
  
\end{figure}
 

Quantum computation follows the circuit model, which is analogous to classical logic circuits. Computation is performed by applying quantum gates (unitary operators) acting on one or more qubits to transform the quantum state. A matrix $U$ is \emph{unitary} if
$
U^\dagger U = U U^\dagger = I
$. As shown in Figure~\ref{fig:gates}, a quantum gate acting on \(N\) qubits
is represented by a \(2^N \times 2^N\) unitary matrix \(U\). It maps an
input state \(\lvert\psi\rangle \in (\mathbb{C}^{2})^{\otimes N} \) to an output state
\(U\lvert\psi\rangle \in (\mathbb{C}^{2})^{\otimes N} \), where \(\mathbb{C}\) denotes
the complex numbers, e.g., one qubit corresponds to $\mathbb{C}^2$ and two qubits to $\mathbb{C}^4$.
A gate set is called universal if it can approximate any unitary operator to arbitrary precision \cite{nielsen2010quantum}. A widely used universal set is the Clifford gates (including the Pauli gates ($X,Y,Z$), the Hadamard gate $H$, the phase gate $S$, the controlled-NOT (CNOT) gate), and the $T$ gate. Figure~\ref{fig:gates} illustrates several of these gates and their matrix representations. A \emph{quantum circuit} is a sequence of quantum gates, and the overall computation corresponds to the composition of their unitary matrices. 

\emph{Circuit depth} is the number of sequential layers of quantum gates in a circuit, representing how many time steps are required to execute it when gates that act on different qubits can be applied in parallel. For example, in Figure~\ref{fig:subcircuits}, the circuit depth is 3. 
We can consider quantum circuits as higher-level operators composed of subcircuits. By abstracting away from hardware and physical implementation details, we can treat such structures modularly. Take the GHZ state preparation circuit in Figure~\ref{fig:subcircuits} for instance. It can serve as a state preparation subcircuit for larger algorithms, i.e., it prepares an input state for other algorithms.


% From a computational perspective, quantum computing uses the quantum circuit model, which is based on gates and circuits analogous to classical logical gates. These gates are operators that act on quantum register states rather than classical bit registers. Quantum registers are composed of quantum bit lines, called qubits. Such quantum states are commonly represented using the dense state vector representation. Consequently, quantum gates are unitary operators that take an input quantum state and output another quantum state. This computation can be mathematically represented using vectors and matrices that live in, and act on, a Hilbert space of dimensionality \(2^N\), where \(N\) is the number of qubits. Like classical bits, qubits are fundamental units of information in quantum logic. Some common gates include the H, S, and CNOT gates; together with the T gate, these form a universal gate set for quantum computation (the Clifford+T gate set), analogous to the role of the NAND gate in classical logic \cite{nielsen2000quantum}. Figure~\ref{fig:gates} displays some of these gates and their operations in matrix notation. 








% % \subsection{Limitations of Existing Benchmarks}
% % \label{ssec:limit}
% % Several quantum circuit benchmarks are widely used in the quantum computing community, including SupermarQ~\cite{tomesh2022supermarqscalablequantumbenchmark}, MQT Bench~\cite{Quetschlich_2023}, and QASM Bench~\cite{li2022qasmbenchlowlevelqasmbenchmark}. While these benchmarks have been valuable for evaluating quantum circuit simulation, they exhibit several limitations.

% % \medskip
% % \emph{Gap 1: No database-native workloads.}  
% % Existing benchmarks primarily distribute quantum circuits as OpenQASM\footnote{\url{https://openqasm.com/}} and  Qiskit QPY\footnote{\url{https://quantum.cloud.ibm.com/docs/en/api/qiskit/qpy}} files. 
% % % These representations are well suited for quantum toolchains, but 
% % They are poorly aligned with database-centric research. In particular, they do not expose simulation workloads in a form that can be directly consumed, queried, or optimized by relational database systems. As a result, studying quantum circuit simulation inside RDBMSs currently requires significant manual translation effort and deep quantum expertise. A benchmark that directly produces simulation workloads in a database-native form, such as SQL queries, would substantially lower the entry barrier for database researchers and enable systematic investigation of query optimization, execution strategies, and storage layouts for simulating quantum computation.

% % \medskip
% % \emph{Gap 2: Lack of structured compositional circuits.}  
% % Real quantum programs are rarely single algorithm instances executed in isolation. Instead, they are typically composed as pipelines of basic subroutines, such as state preparation, algorithmic cores, estimation procedures, and measurement. Take the Harrow–Hassidim–Lloyd (HHL) algorithm~\cite{harrow2009quantum} for example, whose main components are state preparation, Hamiltonian simulation, quantum phase estimation, and measurement. However, existing benchmarks largely treat circuits as independent, fixed templates corresponding to individual algorithms. They do not support systematic composition of building blocks into larger circuits, nor do they model dependencies between subcircuits. Simply sampling circuits independently do not fully represent realistic program structure. In practice, the choice of a subroutine often constrains what follows; for example, a Hamiltonian simulation block is naturally followed by quantum phase estimation to extract spectral information. A benchmark should therefore support conditional, history-dependent circuit generation that reflects such algorithmic structure.

% % \medskip
% % \emph{Gap 3: Limited scalability.}  \rh{add data-centric aspect, add refs for using ML for }
% % Existing benchmarks focus on fixed algorithm templates, thus, their scalability is inherently limited. The number of available circuits grows slowly, and the diversity of circuit instances is constrained. 
% % % This restricts their usefulness for data-driven approaches that require large, diverse training and evaluation corpora, such as learning-based cost modeling or simulator selection. 
% % A benchmark built around compositional generation can scale naturally by varying subroutine combinations.
% % % , repetition counts, and global constraints, enabling orders-of-magnitude growth in the number of generated circuits while maintaining semantic structure.

% % \medskip
% % \emph{Gap 4: Missing the fact that circuits as graphs.}  
% % Quantum circuits are inherently graph-structured objects, where qubits correspond to vertices and multi-qubit gates induce interactions~\cite{Markov+2008}. This structure is central to quantum circuit simulation and compilation. 
% % % particularly for tensor-network–based methods. 
% % Nevertheless, existing benchmarks expose few graph-derived properties and provide limited support for analyzing structural diversity across circuits. A benchmark that explicitly models and stores graph-level features, such as interaction patterns and connectivity, is necessary to enable database-style analysis and to study how circuit structure impacts simulation cost.

% % \medskip
% % % \para{Our design goal} 
% % \para{Summary}
% % Existing benchmarks  fall short of supporting  database-driven and data-centric  research on quantum circuit simulation. Addressing these gaps requires a new benchmark that \emph{supports compositional circuit generation, scales to large and diverse workloads, exposes graph and structural properties, and integrates naturally with database systems}.


% \subsection{Limitations of Existing Benchmarks}
% \label{ssec:limit}

% Several quantum circuit benchmarks are widely used in the quantum computing community, including
% SupermarQ~\cite{tomesh2022supermarqscalablequantumbenchmark}, MQT Bench~\cite{Quetschlich_2023}, and
% QASM Bench~\cite{li2022qasmbenchlowlevelqasmbenchmark}. While these suites are valuable for evaluating
% quantum circuits, they expose limited support for database-driven and data-centric studies of circuit simulation.

% \medskip
% \emph{Gap 1: No database-native workloads.}
% Existing benchmarks primarily distribute circuits as OpenQASM\footnote{\url{https://openqasm.com/}} and
% Qiskit QPY\footnote{\url{https://quantum.cloud.ibm.com/docs/en/api/qiskit/qpy}} artifacts. These formats
% are well suited for quantum toolchains, but they are poorly aligned with database research: they do not expose
% simulation workloads in a form that can be directly executed, inspected, and optimized by an RDBMS.
% As a result, studying circuit simulation inside an RDBMS typically requires substantial manual translation and
% non-trivial quantum expertise.
% \emph{We address this gap by having \sys generate circuits together with an RDBMS-ready SQL workload
% (Section~\ref{ssec:sql}).}

% \medskip
% \emph{Gap 2: Lack of structured compositional circuits.}
% Real quantum programs are rarely single algorithm instances executed in isolation. Instead, they are typically
% assembled as pipelines of subroutines such as state preparation, algorithmic cores, estimation procedures, and
% measurement. For example, the Harrow--Hassidim--Lloyd (HHL) algorithm~\cite{harrow2009quantum} combines
% state preparation, Hamiltonian simulation, quantum phase estimation, and measurement.
% However, existing benchmarks largely treat circuits as independent, fixed templates for standalone algorithms.
% They do not systematically compose building blocks into larger circuits, nor do they model dependencies between
% subcircuits. In practice, the choice of one subroutine often constrains what can follow (e.g., Hamiltonian simulation
% followed by phase estimation).
% \emph{We address this gap in Section~\ref{sssec:seq_select} by modeling conditional, history-dependent template
% composition as a Markov transition process.}

% \medskip
% \emph{Gap 3: Limited scalability for data-centric analysis.}
% Because many existing benchmarks focus on fixed algorithm templates, their scalability is inherently limited:
% the number of instances grows slowly, and structural diversity is constrained. This limits their usefulness for
% data-centric workflows that require large corpora, such as learned routing, cost modeling, and feature-driven analysis.
% \emph{We address this gap with scalable compositional generation and validate scalability experimentally in
% Section~\ref{sec:scalability}; we also release a large dataset of 202k circuits to support large-scale studies.}

% \medskip
% \emph{Gap 4: Limited circuit-as-graph and SQL-level characterization.}
% Quantum circuits are inherently structured objects: qubits correspond to vertices and multi-qubit gates induce
% interactions~\cite{Markov+2008}. This structure is central to simulation cost and to the shape of the induced relational
% workload. Nevertheless, existing benchmarks expose few graph-derived properties and provide limited support for
% analyzing structural diversity across circuit instances.
% \emph{We address this gap by explicitly extracting and storing graph features (Section~\ref{ssec:graphf}) and
% SQL features for the query representation (Section~\ref{ssec:sql}).}

% \medskip
% \para{Summary}
% Overall, existing benchmarks fall short of supporting database-driven and data-centric research on quantum circuit
% simulation. Addressing these gaps requires a benchmark that produces SQL workloads, supports conditional
% composition of subcircuits, scales to large and diverse instances, and exposes graph- and query-level structure for
% systematic analysis.



\subsection{Limitations of Existing Benchmarks}
\label{ssec:limit}

Several quantum circuit benchmarks are widely used in the quantum computing community, including, for example, 
SupermarQ~\cite{tomesh2022supermarqscalablequantumbenchmark}, MQT Bench~\cite{Quetschlich_2023}, and
QASM Bench~\cite{li2022qasmbenchlowlevelqasmbenchmark}. While these suites are valuable for evaluating
quantum circuits, they expose limited support for database-driven and data-centric studies of circuit simulation.

\medskip
% \emph{Gap 1: No database-native workloads.}
\revisionC{\emph{Gap 1: No broad DBMS-executable benchmark workloads.}}
Existing benchmarks primarily distribute circuits as OpenQASM\footnote{\url{https://openqasm.com/}} and
Qiskit QPY\footnote{\url{https://quantum.cloud.ibm.com/docs/en/api/qiskit/qpy}} files. 
These formats are well-suited for quantum toolchains, but \revisionC{they do not expose the induced simulation computation as optimizer-visible relational workloads that can be directly executed, inspected, or optimized by an RDBMS.}
As a result, studying circuit simulation inside an RDBMS typically requires substantial manual translation and
non-trivial quantum expertise.
\revisionC{
The two recent DB-oriented simulation works discussed above \cite{einsteinSQL2023, hai2025quantum} are important first steps within this gap: they show that SQL/RDBMS execution can support quantum circuit simulation. 
However, their workload coverage is limited. \cite{einsteinSQL2023} targets general Einstein summation and evaluates SQL simulation on Google's Sycamore quantum-supremacy circuit, while \cite{hai2025quantum} studies W and GHZ state preparation and QFT circuits. 
Thus, these works motivate DBMS-backed simulation, but do not provide a benchmark-scale collection of diverse circuits with reusable SQL workloads and workload features.
}

We address this gap by having \sys generate \revisionC{general, diverse} circuits, each coming with an RDBMS-ready SQL workload (Section~\ref{sec:benchmark}). 
\revisionC{
This effort is important because it makes the simulation workload visible to database systems: tensor contractions become join-and-aggregate queries whose plans, indexes, materialization choices, and spill behavior can be studied and optimized. 
Consequently, \sys enables evaluating whether DBMS execution can improve memory usage, runtime, and scalability under memory limits, rather than merely translating circuits into another format.
}
% InferQ addresses this remaining gap by enabling systematic evaluation of query optimization, physical design, materialization, and spill-aware execution for quantum circuit simulation workloads.

% \medskip
\emph{Gap 2: Lack of structured compositional circuits.}
Real quantum programs 
% are rarely single algorithm instances executed in isolation. Instead, they 
are typically
assembled as pipelines of subroutines such as state preparation, algorithmic cores, estimation procedures, and
measurement. 
% For example, Harrow--Hassidim--Lloyd (HHL) algorithm~\cite{harrow2009quantum} combines
% state preparation, Hamiltonian simulation, amplitude amplification, and quantum phase estimation.
However, existing benchmarks largely treat circuits as independent, fixed templates for standalone algorithms.
They do not systematically compose building blocks into larger circuits, nor do they model dependencies between
subcircuits. In practice, the choice of one subroutine often constrains what can follow (e.g., Hamiltonian simulation
followed by phase estimation).
We address this gap in Section~\ref{sssec:seq_select} by modeling conditional, history-dependent template
composition as a Markov transition process.

% \medskip
\emph{Gap 3: Limited scalability for data-centric analysis.}
Because many existing benchmarks focus on fixed algorithm templates, their scalability is inherently limited:
the number of instances grows slowly, and structural diversity is constrained. This limits their usefulness for
data-centric workflows that require large corpora, such as applying machine learning models, cost modeling, and feature-driven analysis. 
We address this gap with scalable compositional generation in Section~\ref{sec:benchmark}, and release a large dataset of \num{202975} circuits to support large-scale quantum circuit simulation studies. We also validate scalability experimentally in Appendix~M of our online technical report \cite{inferqtech}.  
 

% \medskip
\emph{Gap 4: Limited circuit-as-graph and SQL-level characterization.}
Quantum circuits are inherently graph-structured: qubits correspond to vertices and multi-qubit gates induce
interactions as edges~\cite{Markov+2008}. This structure is central to simulation cost and to the shape of the induced SQL 
workload. Nevertheless, existing benchmarks expose few graph-derived properties and provide limited support for
analyzing structural diversity across circuits.
We address this gap by explicitly extracting graph features (Section~\ref{ssec:graphf}) and
SQL features for the query representation (Section~\ref{ssec:sql}).

% \medskip
\para{Summary}
Overall, existing benchmarks fall short of supporting database-driven and data-centric research on quantum circuit
simulation. Addressing these gaps requires a benchmark that produces SQL workloads, supports conditional
composition of subcircuits, scales to large and diverse circuits, and exposes graph- and query-level structure for
systematic analysis.


 